% ============================================================
% Automated Financial Ratio Extraction from Annual Reports
% Using Large Language Models: A Consistency Study
% ============================================================
\documentclass[12pt, a4paper]{article}

% ── Packages ────────────────────────────────────────────────
\usepackage[a4paper, margin=2.5cm]{geometry}
\usepackage{booktabs}
\usepackage{array}
\usepackage{multirow}
\usepackage{multicol}
\usepackage{xcolor}
\usepackage[hidelinks]{hyperref}
\usepackage{amsmath}
\usepackage{amssymb}
\usepackage{caption}
\usepackage{subcaption}
\usepackage{longtable}
\usepackage{tabularx}
\usepackage{colortbl}
\usepackage{rotating}
\usepackage{pdflscape}
\usepackage{enumitem}
\usepackage{parskip}
\usepackage{graphicx}
\usepackage{float}
\usepackage{microtype}
\usepackage{makecell}
\usepackage{pifont}   % \ding for check/cross marks
\usepackage{setspace}

% ── Colours ─────────────────────────────────────────────────
\definecolor{consistent}{HTML}{D6EAD6}   % light green
\definecolor{variable}{HTML}{FAD7D7}     % light red
\definecolor{allnull}{HTML}{E8E8E8}      % light grey
\definecolor{hdrblue}{HTML}{2C3E7A}      % dark blue for headers
\definecolor{lightgray}{HTML}{F4F4F4}

% ── Convenience macros ──────────────────────────────────────
\newcommand{\cmark}{\textcolor{green!60!black}{\ding{51}}}  % ✓
\newcommand{\xmark}{\textcolor{red!70!black}{\ding{55}}}    % ✗
\newcommand{\nmark}{\textcolor{gray}{---}}                   % all-null
\newcommand{\metric}[1]{\textit{#1}}
\newcommand{\company}[1]{\textsc{#1}}

% ── Title ────────────────────────────────────────────────────
\title{%
  \textbf{Automated Financial Ratio Extraction from Annual Reports}\\[0.4em]
  \large Using Large Language Models: A Consistency Study
}
\author{Feng Qiu and Alice (Ying) Xia}
\date{Feb 2026}

% ─────────────────────────────────────────────────────────────
\begin{document}
\maketitle
\thispagestyle{empty}

% ── Abstract ────────────────────────────────────────────────
\begin{abstract}
\noindent
We evaluate the consistency and reliability of an automated pipeline that
extracts eight standardised financial ratios from corporate annual-report PDFs
using Google Gemini 2.5~Flash.  The pipeline applies PDF page identification
via fuzzy text matching, targeted page extraction, and a two-call
chain-of-thought prompting strategy.  Each company is processed ten times at
temperature~0.7 to measure extraction stability.  Across nine multinational
companies spanning technology, automotive, luxury goods, and financial
services, we find that \textbf{40.3\%} of all metric--company cells are
consistently extracted (coefficient of variation below 5\%), \textbf{48.6\%}
are variable, and \textbf{11.1\%} are always null.  Net income achieves
perfect reliability (9/9 companies) while cost-to-income and interest coverage
are the least reliable (0\% and 11\% consistent, respectively).  We identify
five root-cause categories for inconsistency and provide concrete
recommendations for each.
\end{abstract}

\vspace{1em}
\tableofcontents
\newpage

% =============================================================
\section{Introduction}
% =============================================================

Extracting structured financial data from unstructured PDF annual reports is a
long-standing challenge in financial natural language processing.  Traditional
approaches rely on hand-crafted rules or proprietary data providers.
Large language models (LLMs) now offer the possibility of zero-shot extraction
directly from scanned or text-based PDFs.  However, LLMs are stochastic:
identical prompts at non-zero temperature may yield different outputs across
runs, raising concerns about reliability in downstream financial analysis.

This paper reports a controlled experiment in which the same extraction prompt
is applied ten times per company to measure how consistently Gemini 2.5~Flash
extracts eight financial metrics from the annual reports of nine global
companies.  The experiment is implemented as a four-step pipeline
(Sections~\ref{sec:pipeline}--\ref{sec:consistency-measure}) and the results
are analysed in detail in Section~\ref{sec:results}.

% =============================================================
\section{Methodology}
\label{sec:methodology}
% =============================================================

\subsection{Pipeline Overview}

The pipeline consists of four sequential steps, as illustrated in
Figure~\ref{fig:pipeline}.  Raw company annual-report PDFs are ingested and
transformed into structured JSON records containing the extracted financial
metrics and their associated consistency statistics.

\begin{figure}[H]
\centering
\fbox{%
\begin{minipage}{0.88\textwidth}
\centering
\vspace{0.5em}
\textbf{Step 1:} Page Identification (\texttt{table\_identifier.py})\\[0.3em]
$\downarrow$\\[0.3em]
\textbf{Step 2:} Page Extraction (\texttt{extract\_pages.py})\\[0.3em]
$\downarrow$\\[0.3em]
\textbf{Step 3:} Gemini Extraction $\times 10$ (\texttt{step\_3\_gemini\_analyzer.py})\\[0.3em]
$\downarrow$\\[0.3em]
\textbf{Step 4:} Consistency Reporting (\texttt{step\_4\_overall\_report.py})
\vspace{0.5em}
\end{minipage}%
}
\caption{Four-step automated financial ratio extraction pipeline.}
\label{fig:pipeline}
\end{figure}

% ─────────────────────────────────────────────────────────────
\subsection{Step 1 — Financial Statement Page Identification}
\label{sec:step1}

The first step identifies which pages of each annual-report PDF contain the
income statement and balance sheet.  This is implemented in
\texttt{table\_identifier.py} using the \texttt{pdfplumber} library.

\subsubsection{Text-Based Table Extraction}

Standard PDFs use either border-based tables (drawn with lines) or
text-based layouts (columns aligned purely by whitespace).  We found that
text-based extraction with explicit column strategies generalises better across
the diverse report formats in our dataset:

\begin{verbatim}
tables = page.extract_tables({
    "vertical_strategy":   "text",
    "horizontal_strategy": "text",
})
\end{verbatim}

\subsubsection{Fuzzy Page Matching}

For each page, the top five text lines are scored against a keyword vocabulary
for income statements (e.g., ``consolidated statements of income'',
``profit and loss account'') and balance sheets (e.g., ``consolidated balance
sheet'', ``statement of financial position'') using the RapidFuzz
\texttt{ratio} function (full-string match):

\begin{equation}
  s_{\text{page}} = \max_{k \in K,\, l \in L_{\text{top5}}}
      \text{fuzz.ratio}(\text{normalise}(l),\, k)
\end{equation}

where $K$ is the keyword set and $L_{\text{top5}}$ is the set of the five
highest-scoring lines on the page.  To handle PDFs where text is concatenated
without spaces (e.g., \texttt{CONSOLIDATEDSTATEMENTOFINCOME}), the score
is taken as the maximum of the original comparison and a
whitespace-stripped comparison:

\begin{equation}
  s = \max\!\bigl(\text{fuzz.ratio}(l,\, k),\;
                  \text{fuzz.ratio}(\text{strip}(l),\, \text{strip}(k))\bigr)
\end{equation}

A minimum title length of 10 characters is enforced to eliminate
short false-positive matches.

\subsubsection{Top-2 Page Filter}

After scoring all pages, a filter retains at most the top-2 pages per
statement type.  A special case applies: if the top-ranked page has a
strictly higher score than all others \emph{and} all remaining pages share
the same score, only the top page is retained.  This avoids including
auxiliary or duplicate statement pages.

% ─────────────────────────────────────────────────────────────
\subsection{Step 2 — Targeted Page Extraction}
\label{sec:step2}

The identified income-statement and balance-sheet page numbers (physical,
1-indexed) are passed to \texttt{extract\_pages.py}, which uses PyMuPDF
(\texttt{fitz}) to extract only those pages into two separate PDF files per
company:

\begin{itemize}[noitemsep]
  \item \texttt{<Company>\_income\_statement.pdf}
  \item \texttt{<Company>\_balance\_sheet.pdf}
\end{itemize}

Providing only the relevant pages --- rather than the entire annual report ---
reduces token consumption and eliminates irrelevant context that could confuse
the model.

% ─────────────────────────────────────────────────────────────
\subsection{Step 3 — LLM-Based Financial Ratio Extraction}
\label{sec:step3}

\subsubsection{Model and Configuration}

Extraction is performed by Google Gemini~2.5~Flash via the
\texttt{google-genai} SDK.  Key parameters:

\begin{center}
\begin{tabular}{ll}
\toprule
\textbf{Parameter} & \textbf{Value} \\
\midrule
Model              & \texttt{gemini-2.5-flash} \\
Temperature        & 0.7 \\
Top-P              & 0.95 \\
Max output tokens  & 16{,}448 \\
Thinking mode      & Disabled (\texttt{thinking\_budget = 0}) \\
Runs per company   & 10 \\
\bottomrule
\end{tabular}
\end{center}

\subsubsection{Two-Call Chain-of-Thought Architecture}

Each run consists of two sequential Gemini API calls.

\paragraph{Call~1 --- Chain-of-thought reasoning.}
The two extracted PDFs (income statement and balance sheet) are first combined
into a single in-memory PDF using PyMuPDF.  This combined PDF, together with
the page position information, is submitted to Gemini with a structured
two-step prompt:

\begin{itemize}[noitemsep]
  \item \textbf{Step~1:} Identify which line items are required from the
        income statement and balance sheet to compute each of the eight target
        metrics.
  \item \textbf{Step~2:} Read the actual values of those line items from the
        attached pages and calculate every metric, showing the formula and
        numbers used.
\end{itemize}

The model produces a free-text chain-of-thought analysis explaining its
reasoning.

\paragraph{Call~2 --- Structured JSON output.}
The chain-of-thought text from Call~1 is passed to a second Gemini call
(without the PDF) with instructions to format the results as a structured JSON
object.  Each metric entry contains:

\begin{verbatim}
{
  "<metric_key>": {
    "value":   <number or null>,
    "unit":    "<e.g. millions USD, ratio, %>",
    "formula": "<formula used>",
    "inputs":  {"<line_item>": <value>, ...},
    "reason":  "<one sentence interpreting this result>"
  }, ...
}
\end{verbatim}

If a required line item is not present in the extracted pages, the model is
instructed to set \texttt{value} to \texttt{null}.

\subsubsection{Target Metrics}

Eight financial metrics are extracted, covering profitability, liquidity, and
leverage dimensions:

\begin{center}
\begin{tabular}{lll}
\toprule
\textbf{Key} & \textbf{Metric} & \textbf{Formula} \\
\midrule
\texttt{net\_income}             & Net Income             & Bottom-line profit \\
\texttt{cost\_to\_income\_ratio} & Cost-to-Income         & Total costs / Revenue \\
\texttt{quick\_ratio}            & Quick Ratio            & (Current assets $-$ Inventory) / Current liabilities \\
\texttt{debt\_to\_equity\_ratio} & Debt / Equity          & Total debt / Shareholders' equity \\
\texttt{debt\_to\_assets\_ratio} & Debt / Assets          & Total debt / Total assets \\
\texttt{debt\_to\_capital\_ratio}& Debt / Capital         & Total debt / (Debt + Equity) \\
\texttt{debt\_to\_ebitda\_ratio} & Debt / EBITDA          & Total debt / EBITDA \\
\texttt{interest\_coverage\_ratio}& Interest Coverage     & EBIT / Interest expense \\
\bottomrule
\end{tabular}
\end{center}

% ─────────────────────────────────────────────────────────────
\subsection{Step 4 — Consistency Measurement}
\label{sec:consistency-measure}

After ten runs, a consistency report is generated for each company.  For each
metric, the ten extracted values are collected (excluding null values) and the
following statistics computed:

\begin{equation}
  \text{CV} = \frac{\sigma}{\lvert\mu\rvert}
\end{equation}

where $\mu$ and $\sigma$ are the mean and population standard deviation of the
non-null values.  A metric--company cell is classified as \textbf{consistent}
if and only if:

\begin{enumerate}[noitemsep]
  \item The null count across ten runs is zero ($n_{\text{null}} = 0$), and
  \item The coefficient of variation is below the threshold
        $\text{CV} < 0.05$ (5\%).
\end{enumerate}

The threshold was chosen to allow for minor rounding differences in the final
decimal place while flagging genuine disagreements in the extracted values.

% =============================================================
\section{Experimental Setup}
\label{sec:setup}
% =============================================================

\subsection{Dataset}

Nine companies were selected to span a broad range of sectors, reporting
currencies, and accounting frameworks:

\begin{center}
\begin{tabular}{llll}
\toprule
\textbf{Company} & \textbf{Year} & \textbf{Sector} & \textbf{Currency} \\
\midrule
Alibaba  & 2024 & Technology / E-commerce & CNY \\
EM       & 2024 & Technology              & USD \\
GM       & 2023 & Automotive              & USD \\
Google   & 2024 & Technology              & USD \\
JPM      & 2024 & Banking                 & USD \\
LVMH     & 2024 & Luxury Goods            & EUR \\
Meta     & 2024 & Technology              & USD \\
Tencent  & 2024 & Technology              & CNY \\
Volkswagen & 2024 & Automotive / Finance  & EUR \\
\bottomrule
\end{tabular}
\end{center}

Each company's annual report was obtained as a multi-hundred-page PDF.  The
pipeline extracted 1--2 income-statement pages and 1--2 balance-sheet pages per
company for input to Gemini.

\subsection{Evaluation Scale}

Each of the 72 cells (9 companies $\times$ 8 metrics) was assigned one of
three outcome labels:

\begin{center}
\begin{tabular}{clp{8cm}}
\toprule
\textbf{Symbol} & \textbf{Label} & \textbf{Criterion} \\
\midrule
\cmark & Consistent & CV $< 5\%$ and $n_{\text{null}} = 0$ \\
\xmark & Variable   & CV $\geq 5\%$ or $n_{\text{null}} > 0$ \\
\nmark & All-null   & All 10 runs returned \texttt{null} \\
\bottomrule
\end{tabular}
\end{center}

% =============================================================
\section{Results}
\label{sec:results}
% =============================================================

\subsection{Overall Headline}

Table~\ref{tab:headline} summarises the aggregate outcomes across all 72
metric--company cells.

\begin{table}[H]
\centering
\caption{Overall extraction outcome across 72 cells (9 companies $\times$ 8 metrics).}
\label{tab:headline}
\begin{tabular}{lccc}
\toprule
\textbf{Outcome}      & \textbf{Count} & \textbf{Percentage} & \textbf{Symbol} \\
\midrule
Consistently extracted & 29 & 40.3\% & \cmark \\
Variable / noisy       & 35 & 48.6\% & \xmark \\
Always null            &  8 & 11.1\% & \nmark \\
\midrule
\textbf{Total}         & 72 & 100\%  & \\
\bottomrule
\end{tabular}
\end{table}

\subsection{Consistency Grid}

Table~\ref{tab:grid} presents the full consistency grid.  Company names are
abbreviated: ALI~=~Alibaba, EM~=~EM, GM~=~GM, GOO~=~Google, JPM~=~JPM,
LVM~=~LVMH, MET~=~Meta, TEN~=~Tencent, VW~=~Volkswagen.
Cells are shaded: \colorbox{consistent}{\strut\ green~} = consistent,
\colorbox{variable}{\strut\ red~} = variable,
\colorbox{allnull}{\strut\ grey~} = all-null.

\begin{table}[H]
\centering
\caption{Consistency grid: \cmark\,=\,consistent (CV$<$5\%),
         \xmark\,=\,variable, \nmark\,=\,all-null.}
\label{tab:grid}
\setlength{\tabcolsep}{5pt}
\begin{tabular}{lcccccccccc}
\toprule
\textbf{Metric} &
  \textbf{ALI} & \textbf{EM} & \textbf{GM} & \textbf{GOO} &
  \textbf{JPM} & \textbf{LVM} & \textbf{MET} & \textbf{TEN} & \textbf{VW} \\
\midrule
Net Income
  & \cellcolor{consistent}\cmark & \cellcolor{consistent}\cmark
  & \cellcolor{consistent}\cmark & \cellcolor{consistent}\cmark
  & \cellcolor{consistent}\cmark & \cellcolor{consistent}\cmark
  & \cellcolor{consistent}\cmark & \cellcolor{consistent}\cmark
  & \cellcolor{consistent}\cmark \\
Cost-to-Income
  & \cellcolor{variable}\xmark & \cellcolor{variable}\xmark
  & \cellcolor{variable}\xmark & \cellcolor{variable}\xmark
  & \cellcolor{variable}\xmark & \cellcolor{variable}\xmark
  & \cellcolor{variable}\xmark & \cellcolor{variable}\xmark
  & \cellcolor{variable}\xmark \\
Quick Ratio
  & \cellcolor{variable}\xmark & \cellcolor{consistent}\cmark
  & \cellcolor{variable}\xmark & \cellcolor{consistent}\cmark
  & \cellcolor{variable}\xmark & \cellcolor{variable}\xmark
  & \cellcolor{consistent}\cmark & \cellcolor{variable}\xmark
  & \cellcolor{allnull}\nmark \\
Debt / Equity
  & \cellcolor{consistent}\cmark & \cellcolor{consistent}\cmark
  & \cellcolor{consistent}\cmark & \cellcolor{consistent}\cmark
  & \cellcolor{variable}\xmark & \cellcolor{variable}\xmark
  & \cellcolor{consistent}\cmark & \cellcolor{consistent}\cmark
  & \cellcolor{allnull}\nmark \\
Debt / Assets
  & \cellcolor{consistent}\cmark & \cellcolor{consistent}\cmark
  & \cellcolor{variable}\xmark & \cellcolor{variable}\xmark
  & \cellcolor{variable}\xmark & \cellcolor{variable}\xmark
  & \cellcolor{consistent}\cmark & \cellcolor{consistent}\cmark
  & \cellcolor{variable}\xmark \\
Debt / Capital
  & \cellcolor{consistent}\cmark & \cellcolor{consistent}\cmark
  & \cellcolor{variable}\xmark & \cellcolor{consistent}\cmark
  & \cellcolor{variable}\xmark & \cellcolor{variable}\xmark
  & \cellcolor{consistent}\cmark & \cellcolor{consistent}\cmark
  & \cellcolor{allnull}\nmark \\
Debt / EBITDA
  & \cellcolor{variable}\xmark & \cellcolor{consistent}\cmark
  & \cellcolor{variable}\xmark & \cellcolor{allnull}\nmark
  & \cellcolor{variable}\xmark & \cellcolor{allnull}\nmark
  & \cellcolor{variable}\xmark & \cellcolor{allnull}\nmark
  & \cellcolor{allnull}\nmark \\
Interest Coverage
  & \cellcolor{variable}\xmark & \cellcolor{variable}\xmark
  & \cellcolor{variable}\xmark & \cellcolor{allnull}\nmark
  & \cellcolor{variable}\xmark & \cellcolor{variable}\xmark
  & \cellcolor{variable}\xmark & \cellcolor{consistent}\cmark
  & \cellcolor{variable}\xmark \\
\midrule
\textbf{Consistent} & 4 & 6 & 2 & 4 & 1 & 1 & 5 & 5 & 1 \\
\textbf{(\% of 8)}  & 50 & 75 & 25 & 50 & 13 & 13 & 63 & 63 & 13 \\
\bottomrule
\end{tabular}
\end{table}

\subsection{Per-Metric Results}

Table~\ref{tab:metric-summary} reports consistency statistics aggregated
across the nine companies for each metric.

\begin{table}[H]
\centering
\caption{Per-metric consistency summary across nine companies.}
\label{tab:metric-summary}
\begin{tabular}{lrrrr}
\toprule
\textbf{Metric} &
  \makecell{\textbf{Consistent}\\\textbf{companies}} &
  \makecell{\textbf{\%}\\\textbf{consistent}} &
  \makecell{\textbf{Always-null}\\\textbf{companies}} &
  \makecell{\textbf{Median}\\\textbf{CV}} \\
\midrule
Net Income             & 9 & 100.0\% & 0 & 0.0000 \\
Debt / Equity          & 6 &  66.7\% & 1 & 0.0334 \\
Debt / Capital         & 5 &  55.6\% & 1 & 0.0445 \\
Debt / Assets          & 4 &  44.4\% & 0 & 0.0380 \\
Quick Ratio            & 3 &  33.3\% & 1 & 0.1228 \\
Cost-to-Income         & 0 &   0.0\% & 0 & 1.9048 \\
Debt / EBITDA          & 1 &  11.1\% & 4 & 0.0088 \\
Interest Coverage      & 1 &  11.1\% & 1 & 0.0683 \\
\bottomrule
\end{tabular}
\end{table}

\subsubsection{Net Income (100\% consistent)}

Net income is the only metric achieving universal consistency.  Eight of nine
companies return an identical value across all ten runs (CV~=~0).  The sole
exception, GM, shows a CV of 0.013 caused by ambiguity between consolidated
net income (\$10{,}127M) and net income attributable to stockholders
(\$9{,}840M); both values are well-defined and the variance is negligible.
Net income is prominently labelled as the final line of every income statement,
making it the easiest target for extraction.

\subsubsection{Debt / Equity (67\% consistent)}

Six of nine companies pass the consistency test.  Failures are structurally
motivated:
\begin{itemize}[noitemsep]
  \item \textbf{LVMH} (CV = 0.291): Bimodal split between 0.33 and 0.59,
        reflecting inconsistent inclusion of operating lease obligations in
        the debt numerator.
  \item \textbf{JPM} (CV = 1.095): Bank liabilities make ``debt'' undefined
        in standard corporate terms; one run returns 9.16$\times$ (includes
        deposits) while others return $\approx$1.32$\times$ (bonds only).
  \item \textbf{Volkswagen}: All 10 runs null --- insufficient balance sheet
        data in extracted pages.
\end{itemize}

\subsubsection{Debt / Capital and Debt / Assets (56\% and 44\% consistent)}

These metrics share similar failure patterns with D/E.  GM additionally
exhibits a unit-ambiguity failure (expressing the ratio as 44.58 instead of
0.4458), inflating CV to $\approx$1.9.  The unit problem is mechanically
identical to the cost-to-income issue and is fixable with a prompt change.

\subsubsection{Quick Ratio (33\% consistent)}

Three companies achieve consistency (EM, Google, Meta).  Failures arise from:
\begin{enumerate}[noitemsep]
  \item \textbf{Definition ambiguity}: Which short-term investments count as
        ``quick assets'' differs by company (Alibaba: two clusters at 1.30
        and 1.79).
  \item \textbf{Structural inapplicability}: JPM (a bank) lacks conventional
        current assets/liabilities; CV~=~0.877, range 0.44--5.78.
  \item \textbf{Data unavailability}: Tencent and Volkswagen return null in
        9/10 and 10/10 runs respectively, indicating the balance sheet pages
        do not present current-asset breakdowns in a parseable format.
\end{enumerate}

\subsubsection{Cost-to-Income Ratio (0\% consistent)}

This is the only metric with \emph{zero} consistent companies.  However, the
underlying extraction is largely correct --- the failure is entirely a
unit-expression problem.  Table~\ref{tab:cti} shows that Gemini extracts
the same underlying ratio value in most runs but occasionally switches between
expressing it as a decimal (e.g., 0.8454) and as a percentage (e.g., 84.54).
This single ambiguity inflates CV values to the range 0.80--2.73.

\begin{table}[H]
\centering
\caption{Cost-to-income ratio: extracted values and CV by company.
         The underlying ratio is consistent; only the unit differs.}
\label{tab:cti}
\begin{tabular}{lccccc}
\toprule
\textbf{Company} &
  \textbf{Decimal form} &
  \textbf{\% form} &
  \makecell{\textbf{Runs as}\\\textbf{decimal}} &
  \makecell{\textbf{Runs as}\\\textbf{\%}} &
  \textbf{CV} \\
\midrule
Alibaba    & 0.845--0.880 & 84.54--88.0  & 7 & 3 & 1.905 \\
EM         & 0.847        & 84.68        & 9 & 1 & 2.725 \\
GM         & 0.946        & 94.59        & 5 & 5 & 1.195 \\
Google     & 0.679        & 67.89        & 9 & 1 & 2.725 \\
JPM        & 0.517        & 51.70        & 5 & 5 & 1.195 \\
LVMH       & 0.769--0.777 & 44.74--77.7  & 4 & 4 & 1.232 \\
Meta       & 0.578        & 57.82        & 9 & 1 & 2.725 \\
Tencent    & 0.471--0.974 & 47.10        & 8 & 1 & 2.666 \\
Volkswagen & 0.817--0.944 & 81.68--81.69 & 6 & 4 & 0.802 \\
\bottomrule
\end{tabular}
\end{table}

\subsubsection{Debt / EBITDA (11\% consistent)}

This metric suffers primarily from data unavailability.  EBITDA requires
depreciation and amortisation (D\&A) figures, which are typically disclosed
in the cash flow statement or in notes --- not on the income statement or
balance sheet pages provided to Gemini.  Only EM (which presents D\&A
data on the extracted pages) achieves consistent extraction across all 10
runs (CV~=~0.009).  Four companies return null in every run (Google, LVMH,
Tencent, Volkswagen); four others extract values in only 1--3 of 10 runs.

\subsubsection{Interest Coverage (11\% consistent)}

Only Tencent passes the threshold (CV~=~0.030).  Failures fall into three
sub-categories:
\begin{enumerate}[noitemsep]
  \item \textbf{Structural inapplicability} (JPM, GM, Volkswagen): Financial
        services divisions make it unclear which ``interest expense'' to use.
        GM shows a bimodal split between $\approx$0.76$\times$ (total
        financing costs) and $\approx$10--12$\times$ (automotive segment
        only), yielding CV~=~0.732.
  \item \textbf{Single-run outlier} (Alibaba, LVMH, EM): Nine of ten runs
        agree; one rogue run picks a different EBIT figure (e.g., EM returns
        87.5$\times$ in one run vs. $\approx$63$\times$ in all others).
  \item \textbf{Data unavailability} (Google): All ten runs null --- Google
        has minimal interest expense and it is not clearly displayed on the
        extracted pages.
\end{enumerate}

% ─────────────────────────────────────────────────────────────
\subsection{Per-Company Results}

Table~\ref{tab:company-detail} provides detailed statistics for each
company across all eight metrics.

\begin{table}[H]
\centering
\caption{Per-company extraction results.
         Mean values are shown with the CV in parentheses.
         \cmark~= consistent, \xmark~= variable, \nmark~= all-null.
         Units: net income in millions (local currency); ratios are dimensionless.}
\label{tab:company-detail}
\setlength{\tabcolsep}{4pt}
\begin{tabular}{lcccccccc}
\toprule
\textbf{Metric} &
  \textbf{ALI} & \textbf{EM} & \textbf{GM} & \textbf{GOO} &
  \textbf{JPM} & \textbf{LVM} & \textbf{MET} & \textbf{TEN} & \textbf{VW} \\
\midrule
\multicolumn{9}{l}{\textit{Net Income (millions, local currency)}} \\
\quad Mean  & 71332 & 36010 & 10041 & 100118 & 58471 & 12550 & 62360 & 196467 & 12394 \\
\quad CV    & 0.000 & 0.000 & 0.013 & 0.000  & 0.000 & 0.000 & 0.000 & 0.000  & 0.000 \\
\quad       & \cmark& \cmark& \cmark& \cmark & \cmark& \cmark& \cmark& \cmark & \cmark\\
\midrule
\multicolumn{9}{l}{\textit{Cost-to-Income Ratio}} \\
\quad Mean  & 17.94 & 9.23  & 38.40 & 7.40  & 20.99 & 28.09 & 6.30  & 5.23  & 49.36 \\
\quad CV    & 1.905 & 2.725 & 1.195 & 2.725 & 1.195 & 1.232 & 2.725 & 2.666 & 0.802 \\
\quad       & \xmark& \xmark& \xmark& \xmark& \xmark& \xmark& \xmark& \xmark& \xmark\\
\midrule
\multicolumn{9}{l}{\textit{Quick Ratio}} \\
\quad Mean  & 1.656 & 1.065 & 0.726 & 1.660 & 1.869 & 0.636 & 2.822 & \nmark& \nmark\\
\quad CV    & 0.123 & 0.000 & 0.229 & 0.000 & 0.877 & 0.155 & 0.000 &  ---  &  ---  \\
\quad       & \xmark& \cmark& \xmark& \cmark& \xmark& \xmark& \cmark& \xmark& \nmark\\
\midrule
\multicolumn{9}{l}{\textit{Debt / Equity}} \\
\quad Mean  & 0.167 & 0.196 & 1.893 & 0.033 & 2.143 & 0.434 & 0.158 & 0.283 & \nmark\\
\quad CV    & 0.049 & 0.011 & 0.001 & 0.032 & 1.095 & 0.291 & 0.004 & 0.033 &  ---  \\
\quad       & \cmark& \cmark& \cmark& \cmark& \xmark& \xmark& \cmark& \cmark& \nmark\\
\midrule
\multicolumn{9}{l}{\textit{Debt / Assets}} \\
\quad Mean  & 0.097 & 0.111 & 9.273 & 0.024 & 2.300 & 0.201 & 0.104 & 0.161 & 0.317\\
\quad CV    & 0.010 & 0.001 & 1.904 & 0.053 & 1.830 & 0.295 & 0.013 & 0.038 & 0.000\\
\quad       & \cmark& \cmark& \xmark& \xmark& \xmark& \xmark& \cmark& \cmark& \xmark\\
\midrule
\multicolumn{9}{l}{\textit{Debt / Capital}} \\
\quad Mean  & 0.144 & 0.164 & 13.61 & 0.032 & 11.56 & 0.298 & 0.137 & 0.221 & \nmark\\
\quad CV    & 0.045 & 0.009 & 1.904 & 0.022 & 1.894 & 0.200 & 0.008 & 0.026 &  ---  \\
\quad       & \cmark& \cmark& \xmark& \cmark& \xmark& \xmark& \cmark& \cmark& \nmark\\
\midrule
\multicolumn{9}{l}{\textit{Debt / EBITDA}} \\
\quad Mean  & 1.206 & 0.564 & 13.09 & \nmark& 2.355 & \nmark& 0.416 & \nmark& \nmark\\
\quad CV    & 0.038 & 0.009 & 0.000 &  ---  & 0.032 &  ---  & 0.000 &  ---  &  ---  \\
\quad Nulls & 7/10  & 0/10  & 9/10  & 10/10 & 8/10  & 10/10 & 9/10  & 10/10 & 10/10 \\
\quad       & \xmark& \cmark& \xmark& \nmark& \xmark& \nmark& \xmark& \nmark& \nmark\\
\midrule
\multicolumn{9}{l}{\textit{Interest Coverage}} \\
\quad Mean  & 14.56 & 65.24 & 6.770 & \nmark& 1.432 & 20.26 & 54.74 & 17.64 & 5.844\\
\quad CV    & 0.059 & 0.114 & 0.732 &  ---  & 0.343 & 0.059 & 0.009 & 0.030 & 0.068\\
\quad Nulls & 1/10  & 0/10  & 0/10  & 10/10 & 1/10  & 0/10  & 7/10  & 0/10  & 0/10 \\
\quad       & \xmark& \xmark& \xmark& \nmark& \xmark& \xmark& \xmark& \cmark& \xmark\\
\bottomrule
\end{tabular}
\end{table}

\subsubsection{EM --- Best Performer (6/8 consistent, 75\%)}

EM achieves the highest consistency score.  Six of eight metrics pass the
threshold, including the rarely-consistent Debt/EBITDA (CV~=~0.009).  The
only failures are cost-to-income (universal unit ambiguity) and interest
coverage (one outlier run: 87.5$\times$ vs. $\approx$63$\times$ in the other
nine runs).  EM's clean income statement layout and the availability of D\&A
figures on the extracted pages explain its strong performance.

\subsubsection{Meta and Tencent (5/8 consistent, 63\%)}

Both companies benefit from straightforward balance sheets with low leverage
and clearly labelled line items.  Meta fails on Debt/EBITDA and Interest
Coverage primarily due to sparse data (minimal interest expense not prominently
displayed).  Tencent's quick ratio is unreliable (9/10 null), suggesting
its balance sheet format does not clearly separate current assets.

\subsubsection{Alibaba and Google (4/8 consistent, 50\%)}

Alibaba's moderate performance reflects classification ambiguity in short-term
investment items for the quick ratio and a near-borderline CV for D/Equity
(0.049) and D/Capital (0.045).  Google achieves perfect extraction for
Net Income, Quick Ratio, D/Equity, and D/Capital but returns null for
Debt/EBITDA and Interest Coverage (Google has near-zero debt and negligible
interest expense).

\subsubsection{GM (2/8 consistent, 25\%)}

GM fails on six metrics.  The captive finance arm (GM Financial) introduces
structural ambiguity into quick ratio, interest coverage, and leverage ratios.
Additionally, GM exhibits the unit-ambiguity failure in Debt/Assets and
Debt/Capital alongside cost-to-income, showing the most severe expression of
this issue (50/50 decimal/percentage split).

\subsubsection{JPM, LVMH, and Volkswagen (1/8 consistent, 13\%)}

All three bottom performers achieve only net income consistency.
\begin{itemize}
  \item \textbf{JPM}: Standard corporate ratios are structurally incompatible
        with a globally systemically important bank.  All seven non-income
        metrics fail, with CVs ranging from 0.031 (Debt/EBITDA, but 8/10
        null) to 1.894 (Debt/Capital).
  \item \textbf{LVMH}: A luxury-goods conglomerate where significant lease
        obligations generate a stable bimodal split in all leverage ratios.
        D/Equity, D/Assets, and D/Capital each oscillate between two
        internally consistent but mutually exclusive debt definitions.
  \item \textbf{Volkswagen}: The highest null rate of any company (4 metrics
        always null, 3 more with 9/10 null).  This is a data availability
        issue: VW's complex automotive~+ financial-services structure, combined
        with German IFRS reporting conventions, results in the required line
        items not being present on the extracted pages.
\end{itemize}

% =============================================================
\section{Root Cause Analysis}
\label{sec:root-causes}
% =============================================================

The inconsistencies observed across 35 variable cells reduce to five distinct
root-cause categories, ordered by severity and ease of remediation.

\subsection{Category~A --- Unit / Scale Ambiguity}

\textbf{Affected:} Cost-to-income (all 9 companies), Debt/Assets and
Debt/Capital for GM.

\textbf{Pattern:} Gemini extracts the correct underlying numerical value but
expresses it inconsistently as either a decimal fraction (e.g., 0.9459) or a
percentage (e.g., 94.59).  At temperature~0.7, the model's output distribution
covers both conventions, leading to catastrophically high CV values despite
the actual information being correct.

\textbf{Fix:} Explicit unit specification in the prompt.  For example:\\
\textit{``Express cost-to-income as a percentage on a 0--100 scale.''} \\
\textit{``Express all ratios as decimals in the range 0--1.''}

This is the single highest-impact fix: it would immediately improve
cost-to-income from 0/9 to a predicted 8--9/9 consistent.

\subsection{Category~B --- Debt Definition Ambiguity}

\textbf{Affected:} LVMH D/E, D/A, D/Capital; JPM D/E, D/A, D/Capital.

\textbf{Pattern:} The balance sheet contains two or more legitimate candidates
for ``total debt'' --- typically a narrow definition (short-term borrowings +
long-term debt) and a broad definition (adding lease obligations, deposits,
or other financial liabilities).  Gemini alternates between them depending on
the chain-of-thought reasoning path taken.  For LVMH, this produces clean
bimodal clusters (e.g., 0.33 and 0.59 for D/E in alternating runs).

\textbf{Fix:} Precisely define the debt numerator in the prompt:
\textit{``Total financial debt = short-term borrowings + current portion of
long-term debt + long-term debt.  Exclude operating lease liabilities, trade
payables, and accrued liabilities.''}

\subsection{Category~C --- Structural Inapplicability}

\textbf{Affected:} Most metrics for JPM; Quick Ratio and Interest Coverage
for GM and Volkswagen.

\textbf{Pattern:} Standard corporate financial ratios were developed for
non-financial industrial companies.  Applying them to banks (JPM) or to
conglomerates with captive finance arms (GM Financial, VW Financial Services)
yields results that are neither wrong nor right --- different analysts would
use fundamentally different denominators.  The model correctly identifies
multiple plausible calculations and picks one per run, producing high variance.

\textbf{Fix:} This requires metric redesign rather than prompt tuning.
For banks: replace D/E, D/A, Quick Ratio, and Interest Coverage with
institution-specific metrics (Tier-1 Capital Ratio, Net Interest Margin,
Loan-to-Deposit Ratio, Return on Assets).  For automotive conglomerates:
specify segment isolation (``use automotive operations only, exclude financial
services subsidiaries'').

\subsection{Category~D --- Data Unavailability}

\textbf{Affected:} Debt/EBITDA (5 companies null or sparse), Interest
Coverage for Google and Volkswagen, Quick Ratio for Tencent and Volkswagen.

\textbf{Pattern:} The required input line items do not appear on the
extracted income statement and balance sheet pages.  EBITDA requires
depreciation and amortisation (D\&A), which is typically disclosed in the
cash flow statement or supplemental notes, not on the income statement.
Similarly, some companies embed interest expense in notes rather than on
the face of the income statement.

\textbf{Fix:} Extend the page extraction to include the first page of the
cash flow statement (to obtain D\&A) and relevant notes.  This is a data
engineering change, not a model change.

\subsection{Category~E --- Single-Run Outlier}

\textbf{Affected:} Alibaba Interest Coverage (1 outlier run), LVMH Interest
Coverage (1 outlier), EM Interest Coverage (1 outlier), GM Net Income (minor).

\textbf{Pattern:} Nine of ten runs agree on a value; one rogue run selects a
slightly different (but legitimate) line item.  These cells are ``almost
consistent'' and would pass the threshold at lower temperature.

\textbf{Fix:} Reduce temperature from 0.7 to 0.3 or 0.0 for production
deployments.  This category represents the easiest win for the five affected
cells.

% =============================================================
\section{Discussion and Recommendations}
\label{sec:discussion}
% =============================================================

\subsection{Metric Reliability Ranking}

Table~\ref{tab:ranking} ranks the eight metrics by their production readiness
based on observed consistency rates and root-cause severity.

\begin{table}[H]
\centering
\caption{Metric reliability ranking and recommended actions.}
\label{tab:ranking}
\begin{tabular}{clcp{5.5cm}}
\toprule
\textbf{Rank} & \textbf{Metric} & \textbf{Consistent} & \textbf{Action} \\
\midrule
1 & Net Income        & 9/9  & Production-ready as-is \\
2 & Debt / Equity     & 6/9  & Specify debt definition; skip JPM \\
3 & Debt / Capital    & 5/9  & Specify unit + debt definition \\
4 & Debt / Assets     & 4/9  & Specify unit + debt definition \\
5 & Quick Ratio       & 3/9  & Specify quick-asset components \\
6 & Debt / EBITDA     & 1/9  & Add cash flow statement pages \\
7 & Interest Coverage & 1/9  & Add interest expense pages; skip banks \\
8 & Cost-to-Income    & 0/9  & Add unit instruction (highest ROI fix) \\
\bottomrule
\end{tabular}
\end{table}

\subsection{Company Suitability for Automated Extraction}

Technology companies with simple, straightforward financial statements (EM,
Meta, Tencent, Google, Alibaba) are well-suited to the current pipeline.
Conglomerates with complex capital structures (LVMH, Volkswagen) require
additional prompt engineering.  Financial institutions (JPM) require
fundamentally different metrics.

\subsection{Prioritised Improvement Roadmap}

\begin{description}[style=nextline]
  \item[\textbf{P1 --- Prompt: unit specification} (High impact, zero code change)]
    Add explicit unit instructions for cost-to-income, Debt/Assets, and
    Debt/Capital.  Expected to fix cost-to-income for all 9 companies and
    Debt/Assets/Capital for GM.  Estimated improvement: $+$13 consistent cells.

  \item[\textbf{P2 --- Prompt: debt definition} (High impact, zero code change)]
    Precisely define the debt numerator.  Expected to stabilise D/E, D/A,
    D/Capital for LVMH and partially for JPM.  Estimated improvement:
    $+$3--6 cells.

  \item[\textbf{P3 --- Data: add cash flow pages} (Medium effort)]
    Include the first page of the operating cash flow section.  Required for
    reliable Debt/EBITDA extraction.  Expected improvement: $+$4--5 cells.

  \item[\textbf{P4 --- Metric redesign for banks} (High effort)]
    Replace all seven non-income metrics for JPM with bank-specific KPIs.
    Expected improvement: $+$6 cells for JPM alone.

  \item[\textbf{P5 --- Temperature reduction} (Zero effort)]
    Reduce temperature from 0.7 to 0.3.  Eliminates Category~E single-run
    outliers.  Estimated improvement: $+$3--5 cells.
\end{description}

% =============================================================
\section{Conclusion}
\label{sec:conclusion}
% =============================================================

This study evaluates the reproducibility of LLM-based financial ratio
extraction across nine global companies and eight metrics over ten independent
runs at temperature~0.7.  The key findings are:

\begin{enumerate}
  \item \textbf{Net income extraction is production-ready} across all sectors
        (100\% consistent).
  \item \textbf{Leverage ratios are moderately reliable} for non-financial
        companies: D/Equity achieves 67\%, D/Capital 56\%, D/Assets 44\%.
  \item \textbf{Cost-to-income fails universally} due to a single, fixable
        unit-ambiguity problem.  The underlying extracted value is correct;
        only its scale fluctuates.
  \item \textbf{Debt/EBITDA and Interest Coverage are data-limited}, not
        model-limited.  Including cash flow statement pages would
        substantially improve both metrics.
  \item \textbf{Sector matters}: Technology companies (EM, Meta, Google,
        Tencent) significantly outperform financial institutions (JPM) and
        complex automotive conglomerates (Volkswagen), which require
        sector-specific metric definitions.
  \item \textbf{Three targeted prompt changes} --- unit specification, debt
        definition, and quick-asset enumeration --- are predicted to raise
        overall consistency from 40\% to approximately 60--70\% without any
        code modifications.
\end{enumerate}

The two-call chain-of-thought architecture proves effective at guiding the
model through multi-step financial calculations, and the pipeline's modular
design allows each identified weakness to be addressed independently.  Future
work should evaluate the effect of lower temperatures on consistency, extend
page coverage to include notes and cash flow statements, and develop
sector-specific extraction profiles for financial institutions.

% ============================================================
\appendix
\section{Raw Extracted Values per Company}
\label{app:raw}
% ============================================================

Tables~\ref{tab:ali-raw}--\ref{tab:vw-raw} list the complete set of extracted
values across all ten runs for each company.

\begin{table}[H]
\centering\small
\caption{Alibaba 2024 --- all extracted values (10 runs).}
\label{tab:ali-raw}
\begin{tabular}{lp{8.5cm}rrc}
\toprule
\textbf{Metric} & \textbf{Values} & \textbf{Mean} & \textbf{CV} & \\
\midrule
Net Income         & 71332 $\times$ 10                                                & 71332 & 0.000 & \cmark \\
Cost-to-Income     & 0.845, 0.845, 0.880, 0.880, 88.0, 0.845, 84.54, 0.845, 0.845, 0.880 & 17.94 & 1.905 & \xmark \\
Quick Ratio        & 1.786, 1.303, 1.786, 1.786, 1.790, 1.446, 1.786, 1.303, 1.786, 1.790 & 1.656 & 0.123 & \xmark \\
Debt / Equity      & 0.173, 0.155, 0.173, 0.173, 0.173, 0.155, 0.173, 0.155, 0.173, 0.170 & 0.167 & 0.049 & \cmark \\
Debt / Assets      & 0.097, 0.097, 0.097, 0.097, 0.097, 0.097, 0.097, 0.097, 0.097, 0.100 & 0.097 & 0.010 & \cmark \\
Debt / Capital     & 0.148, 0.134, 0.148, 0.148, 0.148, 0.134, 0.148, 0.134, 0.148, 0.150 & 0.144 & 0.045 & \cmark \\
Debt / EBITDA      & 1.174, 1.174, 1.270, null $\times$ 7                              & 1.206 & 0.038 & \xmark \\
Interest Coverage  & 16.98, 14.26, 14.26, 14.26, 14.26, 14.26, 14.26, 14.26, 14.26, null & 14.56 & 0.059 & \xmark \\
\bottomrule
\end{tabular}
\end{table}

\begin{table}[H]
\centering\small
\caption{EM 2024 --- all extracted values (10 runs).}
\label{tab:em-raw}
\begin{tabular}{lp{8.5cm}rrc}
\toprule
\textbf{Metric} & \textbf{Values} & \textbf{Mean} & \textbf{CV} & \\
\midrule
Net Income         & 36010 $\times$ 10                                                 & 36010 & 0.000 & \cmark \\
Cost-to-Income     & 0.847 $\times$ 9, 84.68                                           &  9.23 & 2.725 & \xmark \\
Quick Ratio        & 1.065 $\times$ 9, 1.065                                           & 1.065 & 0.000 & \cmark \\
Debt / Equity      & 0.196, 0.203, 0.196, 0.196, 0.196, 0.196, 0.196, 0.196, 0.196, 0.196 & 0.196 & 0.011 & \cmark \\
Debt / Assets      & 0.111 $\times$ 9, 0.111                                           & 0.111 & 0.001 & \cmark \\
Debt / Capital     & 0.164, 0.169, 0.164, 0.164, 0.164, 0.164, 0.164, 0.164, 0.164, 0.164 & 0.164 & 0.009 & \cmark \\
Debt / EBITDA      & 0.570, 0.560, 0.560, 0.570, 0.560, 0.570, 0.560, 0.570, 0.560, 0.560 & 0.564 & 0.009 & \cmark \\
Interest Coverage  & 63.17, 63.17, 63.17, 61.59, 63.17, 62.17, 87.48, 63.17, 63.17, 62.17 & 65.24 & 0.114 & \xmark \\
\bottomrule
\end{tabular}
\end{table}

\begin{table}[H]
\centering\small
\caption{GM 2023 --- all extracted values (10 runs).}
\label{tab:gm-raw}
\begin{tabular}{lp{8.5cm}rrc}
\toprule
\textbf{Metric} & \textbf{Values} & \textbf{Mean} & \textbf{CV} & \\
\midrule
Net Income        & 10127 $\times$ 7, 9840 $\times$ 3                                 & 10041 & 0.013 & \cmark \\
Cost-to-Income    & 94.59, 0.946, 94.59, 0.946, 94.59, 94.59, 0.946, 0.946, 0.946, 0.946 & 38.40 & 1.195 & \xmark \\
Quick Ratio       & 0.830, 0.825, 0.410, 0.825, 0.825, 0.830, 0.825, 0.651, 0.825, 0.411 & 0.726 & 0.229 & \xmark \\
Debt / Equity     & 1.890, 1.894, 1.890, 1.894, 1.894, 1.890, 1.894, 1.894, 1.894, 1.894 & 1.893 & 0.001 & \cmark \\
Debt / Assets     & 0.450, 0.446, 44.58, 0.446, 0.446, 44.58, 0.446, 0.446, 0.446, 0.446 & 9.273 & 1.904 & \xmark \\
Debt / Capital    & 0.650, 0.654, 65.44, 0.654, 0.654, 65.44, 0.654, 0.654, 0.654, 0.654 & 13.61 & 1.904 & \xmark \\
Debt / EBITDA     & 13.09, null $\times$ 9                                             & 13.09 & 0.000 & \xmark \\
Interest Coverage & 0.760, 0.757, 10.21, 0.757, 11.42, 10.21, 0.757, 10.21, 12.42, 10.21 & 6.770 & 0.732 & \xmark \\
\bottomrule
\end{tabular}
\end{table}

\begin{table}[H]
\centering\small
\caption{Google 2024 --- all extracted values (10 runs).}
\label{tab:goo-raw}
\begin{tabular}{lp{8.5cm}rrc}
\toprule
\textbf{Metric} & \textbf{Values} & \textbf{Mean} & \textbf{CV} & \\
\midrule
Net Income         & 100118 $\times$ 10                                                & 100118 & 0.000 & \cmark \\
Cost-to-Income     & 0.679 $\times$ 9, 67.89                                           &   7.40 & 2.725 & \xmark \\
Quick Ratio        & 1.661, 1.661, 1.661, 1.661, 1.660, 1.661, 1.661, 1.660, 1.661, 1.660 &  1.660 & 0.000 & \cmark \\
Debt / Equity      & 0.034 $\times$ 9, 0.030                                           &  0.033 & 0.032 & \cmark \\
Debt / Assets      & 0.024 $\times$ 9, 0.020                                           &  0.024 & 0.053 & \xmark \\
Debt / Capital     & 0.032 $\times$ 9, 0.030                                           &  0.032 & 0.022 & \cmark \\
Debt / EBITDA      & null $\times$ 10                                                  &  ---   &  ---  & \nmark \\
Interest Coverage  & null $\times$ 10                                                  &  ---   &  ---  & \nmark \\
\bottomrule
\end{tabular}
\end{table}

\begin{table}[H]
\centering\small
\caption{JPM 2024 --- all extracted values (10 runs).}
\label{tab:jpm-raw}
\begin{tabular}{lp{8.5cm}rrc}
\toprule
\textbf{Metric} & \textbf{Values} & \textbf{Mean} & \textbf{CV} & \\
\midrule
Net Income         & 58471 $\times$ 10                                                 & 58471 & 0.000 & \cmark \\
Cost-to-Income     & 0.517, 0.517, 51.70, 0.517, 51.70, 0.517, 0.517, 51.69, 51.69, 0.517 & 20.99 & 1.195 & \xmark \\
Quick Ratio        & 0.471, 3.080, 0.580, 1.970, 5.780, 1.970, 1.970, 0.555, 0.443, null & 1.869 & 0.877 & \xmark \\
Debt / Equity      & 1.318, 1.320, 1.320, 1.320, 1.160, 9.160, 1.318, 1.320, 1.877, 1.318 & 2.143 & 1.095 & \xmark \\
Debt / Assets      & 0.114, 0.110, 0.110, 0.114, 10.03, 0.789, 0.114, 11.35, 0.162, 0.114 & 2.300 & 1.830 & \xmark \\
Debt / Capital     & 0.569, 0.570, 0.570, 0.569, 53.80, 0.902, 0.569, 56.85, 0.652, 0.569 & 11.56 & 1.894 & \xmark \\
Debt / EBITDA      & 2.280, 2.430, null $\times$ 8                                     &  2.36 & 0.032 & \xmark \\
Interest Coverage  & 1.741, 0.740, 1.740, 0.740, 1.850, 0.741, 1.850, 1.741, 1.741, null & 1.432 & 0.343 & \xmark \\
\bottomrule
\end{tabular}
\end{table}

\begin{table}[H]
\centering\small
\caption{LVMH 2024 --- all extracted values (10 runs).}
\label{tab:lvm-raw}
\begin{tabular}{lp{8.5cm}rrc}
\toprule
\textbf{Metric} & \textbf{Values} & \textbf{Mean} & \textbf{CV} & \\
\midrule
Net Income        & 12550 $\times$ 10                                                  & 12550 & 0.000 & \cmark \\
Cost-to-Income    & 44.74, 0.777, 0.777, 77.70, 76.92, 0.777, 0.777, 0.777, 0.769, 76.92 & 28.09 & 1.232 & \xmark \\
Quick Ratio       & 0.680, 0.677, 0.706, 0.677, 0.680, 0.455, 0.677, 0.677, 0.427, 0.706 & 0.636 & 0.155 & \xmark \\
Debt / Equity     & 0.330, 0.589, 0.331, 0.589, 0.330, 0.589, 0.331, 0.331, 0.589, 0.331 & 0.434 & 0.291 & \xmark \\
Debt / Assets     & 0.150, 0.273, 0.154, 0.273, 0.150, 0.273, 0.154, 0.154, 0.273, 0.154 & 0.201 & 0.295 & \xmark \\
Debt / Capital    & 0.250, 0.371, 0.249, 0.371, 0.250, 0.371, 0.249, 0.249, 0.371, 0.249 & 0.298 & 0.200 & \xmark \\
Debt / EBITDA     & null $\times$ 10                                                   &  ---  &  ---  & \nmark \\
Interest Coverage & 19.86 $\times$ 9, 23.87                                            & 20.26 & 0.059 & \xmark \\
\bottomrule
\end{tabular}
\end{table}

\begin{table}[H]
\centering\small
\caption{Meta 2024 --- all extracted values (10 runs).}
\label{tab:met-raw}
\begin{tabular}{lp{8.5cm}rrc}
\toprule
\textbf{Metric} & \textbf{Values} & \textbf{Mean} & \textbf{CV} & \\
\midrule
Net Income         & 62360 $\times$ 10                                                 & 62360 & 0.000 & \cmark \\
Cost-to-Income     & 0.578 $\times$ 9, 57.82                                           &  6.30 & 2.725 & \xmark \\
Quick Ratio        & 2.822, 2.822, 2.822, 2.822, 2.820, 2.822, 2.822, 2.820, 2.820, 2.822 & 2.822 & 0.000 & \cmark \\
Debt / Equity      & 0.158 $\times$ 9, 0.160                                           & 0.158 & 0.004 & \cmark \\
Debt / Assets      & 0.104 $\times$ 9, 0.100                                           & 0.104 & 0.013 & \cmark \\
Debt / Capital     & 0.136 $\times$ 9, 0.140                                           & 0.137 & 0.008 & \cmark \\
Debt / EBITDA      & 0.416, null $\times$ 9                                            & 0.416 & 0.000 & \xmark \\
Interest Coverage  & 55.08, 55.08, 54.08, null $\times$ 7                             & 54.74 & 0.009 & \xmark \\
\bottomrule
\end{tabular}
\end{table}

\begin{table}[H]
\centering\small
\caption{Tencent 2024 --- all extracted values (10 runs).}
\label{tab:ten-raw}
\begin{tabular}{lp{8.5cm}rrc}
\toprule
\textbf{Metric} & \textbf{Values} & \textbf{Mean} & \textbf{CV} & \\
\midrule
Net Income         & 196467 $\times$ 10                                                & 196467 & 0.000 & \cmark \\
Cost-to-Income     & 0.974, 0.471, 0.471, 0.471, 47.10, 0.471, 0.471, 0.471, 0.471, 0.974 & 5.23 & 2.666 & \xmark \\
Quick Ratio        & 3.793, null $\times$ 9                                            &  3.793 &  ---  & \xmark \\
Debt / Equity      & 0.273, 0.299, 0.285, 0.263, 0.280, 0.285, 0.285, 0.285, 0.290, 0.290 & 0.283 & 0.033 & \cmark \\
Debt / Assets      & 0.161, 0.164, 0.156, 0.156, 0.160, 0.156, 0.156, 0.156, 0.172, 0.172 & 0.161 & 0.038 & \cmark \\
Debt / Capital     & 0.214, 0.230, 0.222, 0.208, 0.220, 0.222, 0.222, 0.222, 0.225, 0.225 & 0.221 & 0.026 & \cmark \\
Debt / EBITDA      & null $\times$ 10                                                  &  ---   &  ---  & \nmark \\
Interest Coverage  & 17.37, 17.37, 18.70, 17.37, 17.37, 17.37, 17.37, 17.37, 18.70, 17.37  & 17.64 & 0.030 & \cmark \\
\bottomrule
\end{tabular}
\end{table}

\begin{table}[H]
\centering\small
\caption{Volkswagen 2024 --- all extracted values (10 runs).}
\label{tab:vw-raw}
\begin{tabular}{lp{8.5cm}rrc}
\toprule
\textbf{Metric} & \textbf{Values} & \textbf{Mean} & \textbf{CV} & \\
\midrule
Net Income         & 12394 $\times$ 10                                                 & 12394 & 0.000 & \cmark \\
Cost-to-Income     & 81.69, 81.68, 81.68, 81.68, 0.944, 81.68, 0.817, 81.68, 0.817, 0.944 & 49.36 & 0.802 & \xmark \\
Quick Ratio        & null $\times$ 10                                                  &  ---  &  ---  & \nmark \\
Debt / Equity      & null $\times$ 10                                                  &  ---  &  ---  & \nmark \\
Debt / Assets      & 0.317, null $\times$ 9                                            & 0.317 & 0.000 & \xmark \\
Debt / Capital     & null $\times$ 10                                                  &  ---  &  ---  & \nmark \\
Debt / EBITDA      & null $\times$ 10                                                  &  ---  &  ---  & \nmark \\
Interest Coverage  & 6.23, 5.53, 5.17, 5.53, 5.88, 5.88, 5.88, 5.53, 6.58, 6.23     &  5.84 & 0.068 & \xmark \\
\bottomrule
\end{tabular}
\end{table}

\end{document}
